%% =====================================================================
%%  T-23-06  The Three Lies
%%  -------------------------------------------------------------------
%%  One idea per file. Core is mandatory. Slots in canonical order.
%%  Max 700 words. No layout commands. No filler. New source material
%%  for this entry goes in sources/B7/T-23-06.tex -- never by
%%  rewriting this file (docs/RULES.md R0.1).
%%  =====================================================================
%% @kind: topic
%% @id: T-23-06
%% @title: The Three Lies
%% @subtitle: Commission, omission, influence --- the old oath names every lie there is
%% @chapter: c23
%% @order: 60
%% @status: stable
%% @sources: B7
%% @updated: 2026-09-19

\topic{T-23-06}{The Three Lies}{Commission, omission, influence --- the old oath names every lie there is}

\begin{core}
B7's taxonomy, read off the witness oath: every lie ever told is one of three strategies. Commission --- the bald falsehood (\emph{tell the truth}). Omission --- the withheld remainder (\emph{the whole truth}). Influence --- statements that manage your perception of the speaker instead of conveying anything about the issue (\emph{nothing but the truth}). The first is famous, the second familiar, and the third --- the most powerful --- almost nobody catches.
\end{core}
\begin{mechanism}
The strategies differ in cost and visibility, which is why operators climb the ladder. Commission is expensive: it is checkable, and when caught it destroys everything else the speaker says. Omission is comfortable: it says true sentences and lets the listener assemble the wrong picture; the workload is quietly transferred to you. Influence is the master strategy --- the speaker answers a question about facts with a statement about themselves: their devotion, their history, their indignation, their faith. Nothing in it is factually false; the lie lives in the transaction, perception managed instead of information conveyed. Its verbal fingerprints recur across B7's catalogue: the nonanswer preamble (\emph{good question, glad you asked}), the perception qualifier (\emph{to be perfectly honest}), the exclusion qualifier that quietly shrinks the claim (\emph{basically, for the most part}), selective memory (\emph{not that I recall}), the referral to some earlier answered occasion, and above all the convincing statement --- the biography offered where the facts were owed. The detection move is one question: \emph{what did that sentence convey about the issue?} If the honest answer is nothing --- it conveyed warmth, outrage, piety, credentials --- you have your strategy label, and per the window rule it counts toward the cluster only in the question's five seconds.
\end{mechanism}
\begin{conditions}
Strongest under direct questioning, where the speaker owes facts on the record. In ambient social talk the taxonomy blurs --- people volunteer self-portraits constantly, and innocence looks identical to influence; the label needs a question pending. It also needs care with culture and register: enthusiasm about one's record is not automatically influence, and the count of indicators, not the presence of one, decides.
\end{conditions}
\begin{application}
  \item Classify every deflection you receive: commission (check it --- commissions are falsifiable), omission (ask the missing-part question), influence (name what it conveyed and re-ask).
  \item Against omission, ask for the whole: \emph{what else?} The remainder is where the lie lives; the first answer was true.
  \item Against influence, refuse the trade: acknowledge the sentiment once --- \emph{I take you at your word --- the question was X} --- and return to the issue.
  \item Keep the fingerprints in your ear: preamble, qualifier, selective recall, referral. Each is a credit card swipe at the perception till.
  \item When you must answer hard questions yourself, know which strategy you are being invited into --- most reputational questions are fished for influence answers.
\end{application}
\begin{feedback}
  \item Working: you hear the preamble pause --- the half-second where an answer is being swapped for a performance.
  \item Working: your re-asks get shorter, because the taxonomy gave you a neutral, non-accusatory route back to the issue.
  \item Failing: you are fact-checking every sentence. Commission is the rarest strategy; you are fighting the last war.
  \item Failing: everything reads as influence. You have left the question-pending condition; without a question owed, a self-portrait is only talk.
\end{feedback}
\begin{failure}
The taxonomy's abuse is cynicism: a person who files all warmth, all credential and all indignation under influence stops hearing anything but strategy, and becomes the impossible interviewee --- technically undefeated, relationally bankrupt. It also mislabels the honest: some people really are devoted, and say so, with no question pending. The instrument is for questions owed answers.
\end{failure}
\begin{countermeasures}
  \item Audit your own answers for a week: how many were facts, how many performances? The strategy you deploy most is the one you will be slowest to see in others.
  \item In any negotiation, count the influence statements per issue. A rising ratio means the facts are running out.
  \item Teach the oath to your own team: it is a checklist for both directions --- what you may not say, and what you should refuse to hear.
  \item When a public figure answers scandal with biography, run the conveyance test out loud; the mismatch is instructive to watch for (T-23-07).
\end{countermeasures}

\sources{B7}
\seesources
\seealso{T-23-05, T-06-02, T-04-04}
